\documentclass[11pt,oneside]{article}
\usepackage[margin=2cm,bindingoffset=0cm]{geometry}
\usepackage{fontspec}
\setmainfont{Frutiger LT Pro}[
  UprightFont = {*Light},
  BoldFont = {*Bold},
  ItalicFont = {*Light Italic},
  BoldItalicFont = {*Bold Italic}]
\usepackage{hyperref}
\usepackage{soul}
\newfontfamily\frutigerregular{FrutigerLTPro}
\hypersetup{
  colorlinks=true,
  linkcolor=blue,
  urlcolor=blue,
  citecolor=blue}
\renewcommand{\UrlFont}{\fontspec{Frutiger LT Pro}[
  UprightFont = {*Light},
  ItalicFont = {*Light Italic}]}
\usepackage{enumitem}
\usepackage{titlesec}
\setlength{\parindent}{0pt}
\setlength{\parskip}{.5pt}
%\titlespacing*{\<section-command>}{left}{before-sep}{after-sep}
\titlespacing*{\section}{0pt}{7pt}{2pt}
\titlespacing*{\subsection}{0pt}{7pt}{2pt}
\titlespacing*{\subsubsection}{0pt}{7pt}{1pt}
\setlist[itemize]{font=\footnotesize, leftmargin=1.75em, topsep=3pt, partopsep=0pt, itemsep=3pt, parsep=1.5pt, after=\vspace{0pt}}  
\setlist[enumerate]{topsep=3pt, partopsep=0pt, itemsep=3pt, parsep=1.5pt, after=\vspace{0pt}}
\renewcommand{\baselinestretch}{1.1}
\usepackage{fancyhdr}
\usepackage{lastpage}
\pagestyle{fancy}
\fancypagestyle{firstpage}{
  \fancyhf{}
  \renewcommand{\headrulewidth}{0pt}
  \fancyfoot[L]{\today}
  \fancyfoot[R]{Page \thepage\ of \pageref{LastPage}}}
\fancypagestyle{main}{
  \fancyhf{}
  \renewcommand{\headrulewidth}{0.75pt}
  \fancyhead[L]{\textbf{Curriculum Vitae - Alexander Abuabara}}
  \fancyfoot[L]{\today}
  \fancyfoot[R]{Page \thepage\ of \pageref{LastPage}}}

\usepackage{chngcntr}
\counterwithout{subsection}{section} % removes the section dependency
\renewcommand{\thesubsection}{\arabic{subsection}} % shows just the subsection number

\begin{document}
\pagestyle{firstpage}
\raggedright

\centerline{\Large \textbf{Curriculum Vitae - Alexander Abuabara}}

\vspace{12pt}

\setcounter{section}{1}

%\subsection{Personal Information}
\subsection{Profile}
%\section*{Alexander Abuabara}
\begin{itemize}[leftmargin=40pt, label={}]
\item Name: Alexander Abuabara\\
% Nationality: American, Brazilian\\
% Birth: Rio de Janeiro, Brazil (June 1976)\\
% Marital Status: Married, one son born 2010\\
      Email: \href{abuabara@tamu.edu}{abuabara@tamu.edu}\\
      Website: \url{https://abuabara.github.io}\\
      Google Scholar: \url{https://scholar.google.com/citations?hl=en&user=r7vp11AAAAAJ}\\
      ORCID: \url{https://orcid.org/0000-0002-3522-6101}\\
      Phone: +1 979.458.5670\\
      Mail: Office 4013, Emerging Technologies Building\\
            \hspace{0.85cm} 101 Bizzell Street, College Station, TX 77843-3131\\
\end{itemize}

\subsection{Education}
\begin{itemize}[leftmargin=40pt, label={-}]
\item {\frutigerregular Ph.D. in Urban and Regional Science}, Texas A\&M University, College Station (2022)\\
      Dissertation: Applications of Causal Bayesian Networks in Urban Planning\\
      Committee: Walter Gillis Peacock (chair), Zenon Medina-Cetina (chair), Douglas F. Wunneburger, David H. Bierling
%      Committee: Walter Gillis Peacock (chair), Zenon Medina-Cetina (chair), Douglas Frank Wunneburger, David Henry Bierling
%      Committee: Walt G. Peacock (chair), Zenon Medina (chair), Douglas Wunneburger, David Bierling
\item {\frutigerregular Doctoral Studies in Production Engineering}, Federal University of S. Carlos, Brazil (2007-2009)\\
      Advisor: Reinaldo Morabito Neto
%      Advisor: Reinaldo Morabito
\item {\frutigerregular Master of Science in Production Engineering}, Federal University of São Carlos, Brazil (2006)\\
      Thesis: A Mixed Integer Programming Approach Applied to Multi-Period Industrial Planning\\[-2pt]
      {\small (original in Portuguese)}\\
      Committee: Reinaldo Morabito Neto (chair), Horacio Hideki Yanasse, Marcos Nereu Arenales, Denise Sato Yamashita
%      Committee: Reinaldo Morabito (chair), Horacio Yanasse, Marcos Arenales, Denise Yamashita
\item {\frutigerregular Bachelor of Science in Civil Engineering}, University of São Paulo, Brazil (2001)
\end{itemize}

\subsection{Experience}
\subsubsection{Academic Positions}
\begin{itemize}[leftmargin=40pt, label={-}]
\item {\frutigerregular {Instructional Assistant Professor}} (Fall 2025-present)\\
      {Wm Michael Barnes ’64 Department of Industrial and Systems Engineering, Texas A\&M University}\\
      Supervisor: Lewis Ntaimo
\end{itemize}

\begin{itemize}[leftmargin=40pt, label={-}]
\item {\frutigerregular Instructional Assistant Professor} (Fall 2024-Spring 2025)\\
	  {Department of Landscape Architecture and Urban Planning, Texas A\&M University}\\
      Supervisor: Galen D. Newman
\item {\frutigerregular Postdoctoral Research Associate} (Fall 2022-Spring 2024)\\
	  {Department of Landscape Architecture and Urban Planning, Texas A\&M University}\\
      Supervisors: Walter Gillis Peacock, Galen D. Newman
\item {\frutigerregular Graduate Research Lecturer} (Fall 2019-Spring 2022)\\
	  {Department of Landscape Architecture and Urban Planning, Texas A\&M University}\\
      Supervisor: Galen D. Newman, Shannon Van Zandt
\item {\frutigerregular Graduate Research Assistant} (Spring 2015-Spring 2022)\\
	  {Department of Landscape Architecture and Urban Planning, Texas A\&M University}\\
      Supervisor: Walter Gillis Peacock
\end{itemize}

\newpage
\pagestyle{main}

\subsubsection{Industry Experience}
\begin{itemize}[leftmargin=40pt, label={-}]
\item {\frutigerregular Civil Engineer} (2002-2005)\\
      Department of Research, Sucocitrico Cutrale LLC, Araraquara, Brazil\\
      Supervisor: Marcos Neves Penteado Moraes\\
\item {\frutigerregular Civil Engineer Trainee} (1999-2000)\\
      PRO-Salas, University of São Paulo, São Carlos, Brazil\\
      Supervisor: Manoel Rodrigues Alves\\
\end{itemize}

\subsection{Teaching}

\subsubsection{Courses Taught}
\hspace{40pt}{\frutigerregular Industrial and Systems Engineering, Texas A\&M University}\\
\hspace{40pt}({Average Teaching Evaluation 4.55/5})
\vspace{1pt}
\begin{itemize}[leftmargin=40pt, label={-}]
\item {\frutigerregular {DAEN 460 Capstone Senior Design}} (Spring 2026)\\
      Credits 3. 1 Lecture Hour. 6 Lab Hours. Second course in a two-course sequence for the capstone senior design experience; continuation of work on the senior design project in teams; data collection, analysis, application of data engineering methods and tools and development of recommendations considering design constraints, evaluation of alternative design, and application of relevant standards; engagement in oral presentations and creation of the project report, with relevant feedback provided during the semester.
\item {\frutigerregular {DAEN 489 / ISEN 489 Spatial Data Science in Engineering}} (Spring 2026)\\
      Credits 3. 3 Lecture Hours. This course provides engineers with both theoretical foundations and practical skills in spatial databases and geographic information systems (GIS), with emphasis on spatial data modeling, storage, querying, analysis, and visualization. Students will learn to use open source tools to handle real spatial datasets, build spatial databases, perform spatial analysis, and integrate GIS solutions in engineering workflows.
\item {\frutigerregular {DAEN 427 / ISEN 427 Decision and Risk Analysis}} (Fall 2025)\\
      Credits 3. 3 Lecture Hours. Overview of the state of the art in descriptive and prescriptive theories of decision making under uncertainty with emphasis on the ways in which human decisions depart from normative models of rationality; analytical foundations stemming from several disciplines, economics, psychology, management science; application in engineering systems will be considered.
\item {\frutigerregular {DAEN 459 Capstone Senior Design Planning}} (Fall 2025)\\
      Credits 3. 2 Lecture Hours. 3 Lab Hours. First in a two-course sequence for the capstone senior design experience; formation of a senior design team, visitation with the team sponsor, preparation of the groundwork for the project, preparation of the project charter and collection or acquisition of initial set of data; provision of instructions on different aspects of capstone design, including ethics, design constraints, applicable standards, project management, report writing specifications and requirements, and oral and visual presentations.
%\item DAEN 300 Data Engineering Coding Experience I (Fall 2025)
%\item DAEN 301 Data Engineering Coding Experience II (Spring 2026)
\end{itemize}
\vspace{8pt}
\hspace{40pt}{\frutigerregular Landscape Architecture and Urban Planning, Texas A\&M University}\\
\hspace{40pt}({Average Teaching Evaluation 4.59/5})
\begin{itemize}[leftmargin=40pt, label={-}]
\item {\frutigerregular URPN 326 / PLAN 626 Advanced GIS in Urban-Regional Planning} (graduate-level; Spring 2025)\\
      Credits 3. 3 Lecture Hours. Advanced instruction in applications of spatial tools for urban planning, landscape architecture, land development, hazard management, and related problems; GIS applications through review of literature and practice; data quality, uncertainty, the integration of GPS, remote sensing and information technology within the context of urban and regional planning.
\item {\frutigerregular URPN 325 Introduction to GIS in Urban and Regional Planning} (Fall 2024, Spring 2025)\\
      Credits 3. 2 Lecture Hours. 3 Lab Hours. Provides an understanding of GIS fundamentals; basic concepts, principles and functions; essential skills for applying GIS in various fields such as urban planning, landscape architecture, land development, environmental studies, transportation and hazard management; based on learning through class projects.
\item {\frutigerregular URPN 310 Urban Analytical Methods II} (Fall 2024, Spring 2025)\\
      Credits 3. 3 Lecture Hours. Focuses on research conducted by planners, sociologists, anthropologists, political scientists and a variety of applied social scientists; examines variety of procedures employed when conducting research in urban areas; furthers understanding and knowledge of statistical methods employed in social research and elements of geographical analysis.
\item {\frutigerregular URPN 210 Urban Analytical Methods I} (Fall and Spring 2019, 2020, 2021, 2022, 2023, 2024)\\
      Credits 3. 3 Lecture Hours. Study of various analytical techniques used in urban and regional decision making; quantitative approaches to analyze and manipulate data; utilization of statistical packages for data, analysis and communication to enhance urban planning modeling.
\item {\frutigerregular PLAN 604 Planning Methods I} (graduate-level; Fall 2023, 2024)\\
      Credits 3. 3 Lecture Hours. Fundamental concepts and methods used in urban and regional research; qualitative and quantitative research designs; measurement and scaling; sampling; data collection; data file construction; introduction to data analysis and statistical inference.
\item {\frutigerregular URSC 642 Analytic Methods in Landscape and Urban Research II} (graduate-level; Spring 2024)\\
      Credits 3. 3 Lecture Hours. Provides a survey of hands on experiences with advanced techniques and procedures related to conceptual measurement and operational issues, data set development and manipulation and data analysis issues critical for conducting academic research.
\item {\frutigerregular URPN 201 Evolving Cities} (Fall 2021, Spring 2024)\\
      Credits 3. 3 Lecture Hours. Introduction to the history of contemporary urban and regional planning and how the evolving forms of cities and regions pose opportunities and/or challenges for planners; understanding key social, economic, political and technological forces that shape city form and function and its ramification for urban and regional planning.
\item {\frutigerregular FYEX 101 First Year Experience} (Fall 2022)\\
      Credits 0. 0 Lecture Hours. Development of self-efficacy, self-awareness and a sense of purpose; active engagement in the learning environment inside and outside of the classroom; social integration within the university community; also taught at Galveston and Qatar campuses. Must be taken on a satisfactory/unsatisfactory basis.
\end{itemize}

\subsubsection{Guest Lecturer}
\begin{itemize}[leftmargin=40pt, label={-}]
\item {\frutigerregular URSC 631 Foundations of Planning}, presenting \emph{Natural Hazards Research} (Spring 2025)\\
      {Department of Landscape Architecture and Urban Planning, Texas A\&M University}\\
      Instructors: Sinan Zhong, Ji Won Nam
\item {\frutigerregular URPN 361 Urban Issues}, presenting \emph{GIS in Planning Communities} (Fall 2022)\\
      {Department of Landscape Architecture and Urban Planning, Texas A\&M University}\\
      Instructor: Madhavi Dixit
\item {\frutigerregular PLAN 610 Structure and Function of Urban Settlements}, presenting \emph{Urban Sprawl} (Fall 2016)\\
      {Department of Landscape Architecture and Urban Planning, Texas A\&M University}\\
      Instructor: Han Park
\end{itemize}

\subsubsection{Teaching Assistant}
\begin{itemize}[leftmargin=40pt, label={-}]
\item {\frutigerregular URSC 642 Analytic Methods II} (Spring 2017, 2018, 2019)\\
      {Department of Landscape Architecture and Urban Planning, Texas A\&M University}\\
      Instructor: Walt G. Peacock
\item {\frutigerregular URPN 201 Evolving Cities} (Spring 2017)\\
      {Department of Landscape Architecture and Urban Planning, Texas A\&M University}\\
      Instructor: Yu Xiao
\item {\frutigerregular Coding, Introduction to Production Engineering} (Spring 2006)\\
      {Department of Production Engineering, Federal University of São Carlos, Brazil}\\
      Instructor: Reinaldo Morabito, Vitoria Pureza
\item {\frutigerregular Operations Research Applied to Management} (Fall 2005)\\
      {Department of Production Engineering, Federal University of São Carlos, Brazil}\\
      Instructor: Reinaldo Morabito
\item {\frutigerregular Operations Research Applied to Logistics} (Fall 2005)\\
      {Department of Production Engineering, Federal University of São Carlos, Brazil}\\
      Instructor: Reinaldo Morabito
\end{itemize}

\subsection{Research}

\subsubsection{Peer Reviewed Journals}
\begin{enumerate}[leftmargin=40pt]
\item Blanks, J., Abuabara, A., Roberts, A., Semien, J. (2021). Preservation at the Intersections: Patterns of Disproportionate Multihazard Risk and Vulnerability in Louisiana's Historic African American Cemeteries. \emph{Environmental Justice, 14} (1), 1-14.
\item Horney, J.A., Karaye, I.M., Abuabara, A., Gearhart, S., Grabich, S., Perez-Patron, M. (2020). The Impact of Natural Disasters on Suicide in the United States. \emph{Crisis, 42} (5), 328-334.
\item Abuabara, A., Abuabara, A., Tonchuk, C. (2017). Comparative Analysis of Death by Suicide in Brazil and in the United States: descriptive, cross-sectional time series study. \emph{Sao Paulo Medical Journal, 135}, 150-156.
\item Abuabara, A., Morabito, R. (2009). Cutting Optimization of Structural Tubes to Build Agricultural Light Aircrafts. \emph{Annals of Operations Research, 169} (1), 149-165.
\item Abuabara, A., Morabito, R. (2008). Modelos de Programação Inteira Mista para o Planejamento do Corte Unidimensional de Tubos Metálicos na Indústria Aeronáutica Agrícola. \emph{Gestão \& Produção, 15} (3), 605-617.
\end{enumerate}

\subsubsection{Research Projects Collaboration}
\begin{enumerate}[leftmargin=40pt]
\item {\emph{Safeguarding Texas's Inland Rural Heritage: Collaborative Strategies for Resilience Against Natural Hazards}} (Spring 2026), funded by the COA Small Grant Program and the Center for Heritage Conservation, co-sponsorship by the Melbern G. Glasscock Center for Humanities Research, Texas A\&M University. PIs: Michelle A. Meyer, Fabrizio Aimar, Alexander Abuabara, John T. Cooper Jr., Seth H. Jordan.
%\item \emph{Hurricane Evacuation Studies:} (i) Southeastern Texas (2024-2025), (ii) Texas Coastal Bend (2018-2019), and (iii) Texas Valley (2014-2016), funded by the {US Army Corps of Engineers (USACE)} and the {Federal Emergency Management Agency (FEMA)}. Co-PIs at TAMU: Walter G. Peacock, J. Andy Mullins (\emph{in memoriam, Andy passed away during the course of this project, his vision and leadership continue to guide our work}), David H. Bierling, Darrell Borchardt, Douglas F. Wunneburger, Alexander Abuabara.
\item \emph{\href{https://texas-planning-atlas-tamu.hub.arcgis.com}{Hurricane Evacuation Studies}:} 
(i) \href{https://texasatlas.arch.tamu.edu}{{Southeastern Texas (2024-2025)}}, 
(ii) \href{https://texas-planning-atlas-tamu.hub.arcgis.com/apps/337900395c604c3cb2cc1feeb41a9053/explore}{Texas Coastal Bend (2018-2019)}, and 
(iii) \href{https://texas-planning-atlas-tamu.hub.arcgis.com/apps/047af4b2415b442e9bb10108ab1d3eb4/explore}{Texas Valley (2014-2016)}, funded by the {US Army Corps of Engineers (USACE)} and the {Federal Emergency Management Agency (FEMA)}.  Co-PIs at TAMU: Walt G. Peacock, J. Andy Mullins (\emph{in memoriam, Andy passed away during the course of this project, his vision and leadership continue to guide our work}), David H. Bierling, Darrell W. Borchardt, Douglas F. Wunneburger, Alexander Abuabara.
\item \emph{Risk-Based Community Resilience Planning} (2022-2024), funded by Colorado State University from the National Institute of Standards and Technology (NIST). Directors: John W. van de Lindt (CSU), Jamie Brown Kruse (CSU); Co-PIs at TAMU: Walt G. Peacock, Shannon S. Van Zandt, Nathanael P. Rosenheim, Michelle A. Meyer, Maria Koliou.
\item \emph{The Adoption and Utilization of Hazard Mitigation Practices by Jurisdictions along Gulf and Atlantic Coasts} (2016-2018), funded by the National Science Foundation. Co-PIs: Walt G. Peacock, Shannon S. Van Zandt, Himanshu Grover (SUNY).
\end{enumerate}

\subsubsection{Reports}
\begin{enumerate}[leftmargin=40pt]
\item Bierling, D.H., Wunneburger, D.F., Peacock, W.G., Borchardt, D.W., Abuabara, A. (2025). {Southeast Texas Hurricane Evacuation Study: Evacuation and Sheltering Behavior Report}. Texas A\&M Hazard Reduction \& Recovery Center and the Texas A\&M Transportation Institute. College Station, TX.
\item Lee, R.J., Day, W., Abuabara, A., Newman, G.D., Peacock, W.G. (2023). The Effect of Flood Buyouts on Nearby Property Values. Working paper for the \emph{Lincoln Institute of Land Policy}.
\item Rosenheim, N., Lane, G., Katare, A., Watson, M., Williams, A., Peacock, W.G., Abuabara, A., Perez, M., Sullivan, E., Kastor, H. (2021). Food Retailer Data, in \emph{Food Access Impact Survey for Harris County and Southeast Texas after Hurricane Harvey in 2017}.
\item Peacock, W.G., Abuabara, A., Wunneburger, D.F., Bierling, D.H., Mullins, J.A., Borchardt, D.W. (2020). Coastal Bend Hurricane Evacuation Study: Evacuation Zone Development Report. Texas A\&M Hazard Reduction \& Recovery Center and the Texas A\&M Transportation Institute. College Station, TX.
\item Peacock, W.G., Wunneburger, D.F., Abuabara, A., Bierling, D.H., Mullins, J.A., Borchardt, D.W. (2020). Coastal Bend Hurricane Evacuation Study: Vulnerability Analysis Report. Texas A\&M Hazard Reduction \& Recovery Center and the Texas A\&M Transportation Institute. College Station, TX.
\item Bierling, D.H., Lindell, M., Peacock, W.G., Abuabara, A., Moore, R., Wunneburger, D.F., Mullins, J. A., Borchardt, D.W. (2020). Coastal Bend Hurricane Evacuation Study: Hurricane Harvey Evacuation Behavior Survey Outcomes and Findings. Texas A\&M Hazard Reduction \& Recovery Center and the Texas A\&M Transportation Institute. College Station, TX.
\item Mullins, J.A., Borchardt, D.W., Bierling, D.H., Peacock, W.G., Wunneburger, D.F., Abuabara, A. (2020). Coastal Bend Hurricane Evacuation Study: Transportation Analysis Report. Texas A\&M Hazard Reduction \& Recovery Center and the Texas A\&M Transportation Institute. College Station, TX.
\item Peacock, W.G., Wunneburger, D.F., Abuabara, A., Park, H., Van Zandt, S. (2016). Rio Grande Valley Hurricane Evacuation Study: Vulnerability Analysis Report. Texas A\&M Hazard Reduction \& Recovery Center and the Texas A\&M Transportation Institute. College Station, TX.
\end{enumerate}

\subsubsection{Data}
\begin{enumerate}[leftmargin=40pt]
\item Lane, G., Rosenheim, N., Williams, A., Watson, M., Abuabara, A., Maher, D., Rathburn, A., Peacock, W.G. (2024). Food Aid Agency Data, in Food Access Impact Survey for Harris County and Southeast Texas after Hurricane Harvey in 2017. DesignSafe-CI.
\end{enumerate}

\subsubsection{Conferences with Peer Reviewed Abstracts}
{\small (underlined name indicates the presenter)}
\begin{enumerate}[leftmargin=40pt]
\item \underline{Wunneburger, D.F.}, Peacock. W.G., Abuabara, A., Bierling, D.H., Borchardt, D.W. Risk-Based Hurricane Evacuation Zone Planning for the Texas Gulf Coast. \emph{2025 Researchers Meeting, organized by the International Sociological Association Research Committee on Disasters (ISA-RC39)}. Broomfield, CO (Jul 2025)
\item \underline{Stanley, M.}, Peacock, W.G., Abuabara, A., Day, W., Wade, H. Procedural and Distributive Equity in Disaster Housing Recovery: FEMA's IA Program for Unmet Housing Needs After Hurricane Ike \emph{50th Annual Natural Hazards Research and Applications Workshop}. Broomfield, CO (Jul 2025)
\item \underline{Day, W.}, Abuabara, A., Peacock, W.G., Wade, H. FEMA's Financial Housing Assistance to Owners and Renters After Hurricane Harvey. \emph{49th Annual Natural Hazards Research and Applications Workshop}. Broomfield, CO (Jul 2024)
\item \underline{Blanks, J.}, Abuabara, A. Did You Check on the Cemetery? A Geospatial Risk Assessment of Cemeteries in Southeast Louisiana. \emph{PRIMR Conference at PVAMU}. Prairie View, TX (Mar 2023)
\item \underline{Abuabara, A.}, Day, W., Peacock, W.G., Wade, H. An Analysis of FEMA's Individual Assistance After Hurricane Harvey. \emph{Natural Hazards Research Summit}. Washington DC (Oct 2022)
\item \underline{Abuabara, A.}, Peacock, W.G., Bierling, D.H., Lindell, M. A Statistical Analysis of Household Evacuation in Response to Hurricane Harvey. \emph{Society for Risk Analysis Annual Meeting}. Austin, TX (Dec 2020)
\item \underline{Abuabara, A.}, Medina-Cetina, Z., Peacock, W.G. Application of Bayesian Networks to Hurricane Risk Assessment. \emph{Society for Risk Analysis Annual Meeting}. Arlington, VA (Dec 2019)
\item \underline{Horney, J.A.}, Karaye, I.M., Gearhart, S., Abuabara, A., Perez-Patron, M. Impact of Natural Disasters on Suicide Risk in the United States. \emph{American Public Health Association Annual Meeting}. Philadelphia, PA (Nov 2019)
\item \underline{Abuabara, A.}, Peacock, W.G. Importance of Mobile Homes to Determine Natural Hazard Evacuation Boundaries. \emph{46th Annual Natural Hazards Research and Applications Workshop}. Broomfield, CO (Jul 2017)
\item \underline{Abuabara, A.}, Considerations for Modeling Natural Hazard Evacuation of Mobile Neighborhoods and Mobile Parks in the Texas Valley Region. \emph{42nd Urban Affairs Assoc. Conference}. San Diego, CA (Mar 2016)
\item Abuabara, A.,\underline{ Morabito, R.} Modelo de Otimização de Corte Unidimensional Aplicado a Fabricação de Aeronaves Agrícolas, \emph{Annals of XXXIX Simpósio Brasileiro de Pesquisa Operacional}. Fortaleza, Brazil (Aug 2007)
\item \underline{Morabito, R.}, Abuabara, A. A One-dimensional Cutting Stock Problem in the Production of Agricultural Light Aircrafts”, \emph{XXII European Conference on Operational Research}, Prague, Czech Republic (Jul 2007)
\item \underline{Abuabara, A.}, Morabito, R. Pipe Cutting Optimization: Application to the Aeronautics Industry. \emph{Latin-Iberoamerican Conference on Operations Research} CLAIO. Montevideo, Uruguay (Nov 2006)
\end{enumerate}

\subsubsection{Workshops, Presentations, Invited Talks, Seminars \& Panel Participation}
\begin{enumerate}[leftmargin=40pt]
\item Designing for Impact: Faculty Perspectives on High-Impact and Reflective Learning (Panelist). Texas A\&M Center for Teaching Excellence. College Station, TX (Jan 27, 2026). Organizer: Nate Poling
\item {Introduction to QGIS. Texas A\&M GIS Day.} College Station, TX (2018, 2019, 2020, 2021, 2025)\\
      Organizer: Daniel Goldberg et al.
\item Problem Structuring, Data, and My Teaching Journey, and Why My Students Still Ask for Extra Credit. ISEN Spring Seminar. College Station, TX (2025)
\item {Risk-based Evacuation Zone Planning for the Texas Gulf Coast.} CoA Research Symposium. College Station, TX (Fall 2025)
\item Zone Planning for Coastal Texas. CoA Research Symposium. College Station, TX (2024) % Hurricane Evacuation Studies: 
\item Application of Bayesian Networks to Hurricane Risk Assessment. CoA Research Symposium. College Station, TX (Fall 2024)
\item Modeling Individual Assistance to Households ​in Texas After Hurricane Harvey. CoA Research Symposium. College Station, TX (Fall 2024)
\item Hurricane Evacuation Studies: Zone Planning for Coastal Texas. CoA Research Symposium. College Station, TX (Fall 2024)
\item Modeling Individual Assistance to Households ​in Texas After Hurricane Harvey. CoA Research Symposium. College Station, TX (Fall 2023)
\item A Web-Based GIS Platform and Applications for Evacuation, Preparedness, Response, Mitigation, and Recovery Planning: The Texas Hazard Resilience Atlases. Texas Emergency Management Conference. College Station and Dallas, TX (Spring 2023)
\item FEMA Individual Assistance Program for Housing. IN-CORE Summer Research Meeting. Fort Collins, CO (Spring 2022)
\item Spatial Analysis in R. Texas A\&M GIS Day. College Station, TX (2018, 2019, 2020, 2021)\\
      Organizer: Daniel Goldberg et al.
\item Hurricane Evacuation Studies. Texas A\&M GEOSAT Seminar. College Station, TX (Fall 2018)

\end{enumerate}

\subsection{Service}
\begin{enumerate}[leftmargin=40pt]
\item {Undergraduate Affairs Committee (Fall 2025-present)}\\
      {Department of Industrial and Systems Engineering, Texas A\&M University}\\
      Chair: Amarnath Banerjee
\item Faculty Senator (Fall 2024-Summer 2025)\\
      {College of Architecture, Texas A\&M University}\\
      COA Faculty Senator Caucus Leader: Shelley D. Holliday\\
      Faculty Senate Speakers: Angie Hill and Andrew Klein
\item Faculty Member: Master of Urban Planning (MUP) reaccreditation, strategic planning (2024-2025)
      {Department of Landscape Architecture and Urban Planning, Texas A\&M University}\\
      Coordinator: Ivis Garcia Zambrana
\item Faculty Member: Bachelor of Urban/Regional Planning accreditation, strategic planning (2024-2025)\\
      {Department of Landscape Architecture and Urban Planning, Texas A\&M University}\\
      Coordinator: Justin Golbabai
\end{enumerate}

\subsection{Students}
\subsubsection{Current Students}
\hspace{40pt}{\frutigerregular External Committee Member}\\
\begin{enumerate}[leftmargin=40pt]
\item Adam Mann, Master of Urban Planning (Expected Spring 2026)
\item Aaron Garcia, Master of Urban Planning (Expected Spring 2026)
%\item Jen-Wei Huang, Master of Urban Planning (Expected Spring 2026)
\end{enumerate}

\subsubsection{Capstone Faculty Advisor}
\hspace{40pt}{\frutigerregular Industrial \& Systems Engineering, Texas A\&M University}\\
\begin{enumerate}[leftmargin=40pt]
\item Jan Barrett, Eric Gomez, Annemarie Hoffzimmer, Thomas Lim, Isha Pokerna, Gabrielle Wibner.
      Pedestrian Conflict at Railroad Crossings. Sponsor: Cyndi Lee, Program Aide for the East End District, Houston, TX.
      \emph{3rd place in major of the 2025 COE Project Showcase} (Fall 2025)
\end{enumerate}

\subsection{Honors, Fellowships, Awards}

\subsubsection{Professional Development}
\begin{enumerate}[leftmargin=40pt]
\item {Certificate in Promoting Active Learning,} Association of College and University Educators (ACUE, Fall 2025)
\item Faculty Affairs Third-Party Conflict Resolution Mediator, Texas A\&M University (Fall 2024)\\
      Associate Provost for Faculty Affairs: Heather Wilkinson
\item Graduate Graduate Mentoring Academy Fellow, Texas A\&M University (Fall 2024)\\
      Assistant Director of Mentoring \& Graduate Education: Clinton A. Patterson
\end{enumerate}

\subsubsection{Professional Affiliations and Memberships}
\begin{enumerate}[leftmargin=40pt]
\item {Faculty fellow for the Transforming Teaching through Engaged and Active Learning (TEAL) Academy} (2025-present)
\item Faculty fellow, Center for Geospatial Sciences, Appl. and Tech., Texas A\&M University (2020-present)
\item Faculty fellow, Stochastic Geomechanics Laboratory, Texas A\&M University (2018-present)
\item Faculty fellow, Hazard Reduction and Recovery Center, Texas A\&M University (2015-present)
\item Society for Risk Analysis (2018-present)
\item Tau Sigma Delta Honor Society (2015-present)
\item Regional Council of Engineering and Agronomy of São Paulo, Brazil (2002-2005)
\end{enumerate}

\subsubsection{Honors and Fellowships}
\begin{enumerate}[leftmargin=40pt]
\item {Travel Award for Natural Hazards Research Summit}, National Science Foundation, Washington DC (Oct 2022)
\item {Summer Institute for Early-Career Faculty Award}, National Science Foundation Natural Hazards Engineering Research Infrastructure (NHERI), San Antonio, TX (Jun 2022)
\item {Open-Source GIS Developers Summer School}, GeoSTAT Institute of Botany of the Czech Academy of Sciences (IBOT), Prague, Czech (Aug 2018)
\item {Coastal Flood Risk Reduction Research}, National Science Foundation Partnerships for International Research and Education (PIRE), Amsterdam, Netherlands (May 2018)
\item {Full Tuition Assistantship}, Texas A\&M University (2015-2022)
\item {Graduate Program Scholarship}, Department of Landscape Architecture and Urban Planning, Texas A\&M University (2014-2015)
\item {Graduate Program Scholarship}, Brazilian Federal Agency for Support and Evaluation of Graduate Education (CAPES), Department of Production Engineering, Federal University of São Carlos, Brazil (2005-2006)
\item {REHIDRO Monitoring Work, Undergraduate Scholarship}, University of São Paulo, Brazil (1997-1998)
\item {CISC Monitoring Work, Undergraduate Scholarship}, University of São Paulo, Brazil (1994-1997)
\end{enumerate}

\end{document}